
\documentclass{article}
\usepackage{colortbl}
\usepackage{makecell}
\usepackage{multirow}
\usepackage{supertabular}

\begin{document}

\newcounter{utterance}

\centering \large Interaction Transcript for game `clean\_up', experiment `0\_easy\_3obj\_tr', episode 0 with gemma{-}4{-}26b{-}a4b{-}it.
\vspace{24pt}

{ \footnotesize  \setcounter{utterance}{1}
\setlength{\tabcolsep}{0pt}
\begin{supertabular}{c@{$\;$}|p{.15\linewidth}@{}p{.15\linewidth}p{.15\linewidth}p{.15\linewidth}p{.15\linewidth}p{.15\linewidth}}
    \# & $\;$A & \multicolumn{4}{c}{Game Master} & $\;\:$B\\
    \hline

    \theutterance \stepcounter{utterance}  
    & & \multicolumn{4}{p{0.6\linewidth}}{
        \cellcolor[rgb]{0.9,0.9,0.9}{
            \makecell[{{p{\linewidth}}}]{
                \texttt{\tiny{[P1$\langle$GM]}}
                \texttt{Ben senin oyun yöneticinim ve aşağıdaki ızgarayı oyun tahtası olarak kullandığınız işbirlikçi bir oyun oynuyorsunuz.} \\
\texttt{```} \\
\texttt{ 1234567} \\
\texttt{╔═╤═══╤═╗} \\
\texttt{║◌│◌◌◌│◌║ 1} \\
\texttt{╟─┼──┬┘◌║ 2} \\
\texttt{║◌│◌◌│◌P║ 3} \\
\texttt{╟─┴──┘◌◌║ 4} \\
\texttt{║◌L◌◌◌◌◌║ 5} \\
\texttt{║◌C◌◌◌◌◌║ 6} \\
\texttt{║◌◌◌◌◌◌◌║ 7} \\
\texttt{╚═══════╝} \\
\\ 
\texttt{```} \\
\\ 
\texttt{* Üst kenarda x koordinatları sağa doğru artar, sağ kenarda ise y koordinatları aşağı doğru artar.} \\
\texttt{* Aşağıdaki nesneler oyun tahtasına rastgele yerleştirilmiştir: 'C', 'L', 'P'.} \\
\\ 
\\ 
\texttt{Diğer oyuncu, nesnelerin farklı rastgele konumlara yerleştirildiği oyun tahtasının bir varyasyonunu görür. Sen diğer oyuncunun tahtasını göremezsin ve onlar da seninkini göremez.} \\
\texttt{**Amaç:**} \\
\\ 
\texttt{Her iki oyuncunun da, kendi oyun tahtalarındaki nesneleri hareket ettirerek aynı türden nesnelerin aynı koordinatlara gelmesini sağlaması gerekir. Ortak bir hedef düzen üzerinde anlaşmak için diğer oyuncuyla iletişim kurmalısınız.} \\
\\ 
\\ 
\texttt{**Kurallar:**} \\
\\ 
\texttt{* Her turda, aşağıdaki iki komuttan yalnızca birini gönderebilirsin:} \\
\texttt{ 1. `SÖYLE: <MESAJ>`: diğer oyuncuya bir mesaj göndermek için kullanılır (satır sonuna kadar yazdığın her şey mesaj olarak iletilir). Ben bu mesajı partnerine ileteceğim.} \\
\texttt{2. `HAMLE: <NESNE>, (<X>, <Y>)`: bir nesneyi yeni bir konuma taşımak için kullanılır. Burada <X> sütunu, <Y> ise satırı ifade eder. Yaptığın hamlenin başarılı olup olmadığını sana söyleyeceğim.} \\
\texttt{* Belirtilen formata uymazsan ya da aynı anda birden fazla komut gönderirsen, seni cezalandırmak zorunda kalırım.} \\
\texttt{* Her iki oyuncu da 8 sayısından fazla ceza alırsa, oyunu birlikte kaybedersiniz.} \\
\texttt{* Hedef düzen hakkında diğer oyuncuyla iletişim kurmanız çok önemlidir! Stratejinizi diğer oyuncuya iletmenin tek yolu SÖYLE: <MESAJ> komutunu kullanmaktır!} \\
\\ 
\\ 
\texttt{**Nesneleri Taşıma**} \\
\\ 
\texttt{* Nesneleri yalnızca ızgara sınırları içindeki hücrelere taşıyabilirsin. Hedef hücre boş olmalıdır; yani yalnızca ‘◌’ sembolünü içermelidir.} \\
\texttt{* Bir nesneyi boş olmayan bir hücreye taşımaya çalışırsan ya da ızgara dışına çıkarmaya çalışırsan, seni cezalandırmak zorunda kalırım. Ardından tekrar deneme hakkın olur.} \\
\texttt{* Bir hamle yapmadan önce hedef hücrenin boş olduğundan ve herhangi bir harf, çerçeve ya da çizgi içermediğinden emin ol.} \\
\\ 
\\ 
\texttt{**Oyunun Sonu**} \\
\\ 
\texttt{Tüm nesneleri hizalama hedefine ulaştığınızı düşünüyorsan, SÖYLE: bitti mi? göndererek diğer oyuncudan oyunu bitirmesini isteyebilirsin. Eğer diğer oyuncu senden oyunu bitirmeni isterse ve sen SÖYLE: bitti! ile yanıt verirsen, oyun sona erer.} \\
\\ 
\texttt{Oyun, 12 tur içinde sona ererse her iki oyuncu da kazanır. Bir tur, iki oyuncunun da geçerli bir komut göndermesi olarak tanımlanır.} \\
\\ 
\texttt{**Puanlama:**} \\
\\ 
\texttt{Aynı nesneler her iki oyun tahtasında ne kadar yakın konumlanırsa, o kadar fazla puan kazanırsınız. Cezalar puanınızı düşürür. Rekoru kırabilir misiniz?} \\
\\ 
\\ 
\texttt{Oyuna başlamak için diğer oyuncuya bir mesaj gönder!} \\
            }
        }
    }
    & \\ \\

    \theutterance \stepcounter{utterance}  
    & \multicolumn{4}{p{0.6\linewidth}}{
        \cellcolor[rgb]{0.9,0.9,0.9}{
            \makecell[{{p{\linewidth}}}]{
                \texttt{\tiny{[P1$\rangle$GM]}}
                \texttt{SÖYLE: Selam! Seninle işbirliği yapmaya hazırım. Ortak bir hedef düzen belirlemek için önce nesnelerin konumlarını paylaşalım mı? Ben kendi tahtamdaki nesnelerin yerlerini sana söyleyeceğim, sen de kendi tahtandakileri söyle. Sonra ortak bir hedef üzerinde anlaşırız.} \\
\\ 
\texttt{Benim tahtamdaki konumlar:} \\
\texttt{C: (2, 6)} \\
\texttt{L: (2, 5)} \\
\texttt{P: (7, 3)} \\
            }
        }
    }
    & & \\ \\

    \theutterance \stepcounter{utterance}  
    & & & \multicolumn{2}{p{0.3\linewidth}}{
        \cellcolor[rgb]{0.9,0.9,0.9}{
            \makecell[{{p{\linewidth}}}]{
                \texttt{\tiny{[GM$|$GM]}}
                \texttt{SÖYLE: Selam! Seninle işbirliği yapmaya hazırım. Ortak bir hedef düzen belirlemek için önce nesnelerin konumlarını paylaşalım mı? Ben kendi tahtamdaki nesnelerin yerlerini sana söyleyeceğim, sen de kendi tahtandakileri söyle. Sonra ortak bir hedef üzerinde anlaşırız.} \\
\\ 
\texttt{Benim tahtamdaki konumlar:} \\
\texttt{C: (2, 6)} \\
\texttt{L: (2, 5)} \\
\texttt{P: (7, 3)} \\
            }
        }
    }
    & & \\ \\

    \theutterance \stepcounter{utterance}  
    & & & \multicolumn{2}{p{0.3\linewidth}}{
        \cellcolor[rgb]{0.9,0.9,0.9}{
            \makecell[{{p{\linewidth}}}]{
                \texttt{\tiny{[GM$|$GM]}}
                \texttt{Invalid format: tail is not empty: '} \\
\\ 
\texttt{Benim tahtamdaki konumlar:} \\
\texttt{C: (2, 6)} \\
\texttt{L: (2, 5)} \\
\texttt{P: (7, 3)} \\
            }
        }
    }
    & & \\ \\

    \theutterance \stepcounter{utterance}  
    & & \multicolumn{4}{p{0.6\linewidth}}{
        \cellcolor[rgb]{0.9,0.9,0.9}{
            \makecell[{{p{\linewidth}}}]{
                \texttt{\tiny{[P1$\langle$GM]}}
                \texttt{Ceza: Mesajında komuttan sonra hiçbir şey bulunmamalı!} \\
\\ 
\texttt{Yanıtının yalnızca SÖYLE: <MESAJ> ya da HAMLE: <NESNE>, (<X>, <Y>) komutlarından birini içerdiğinden emin ol; bunun dışında hiçbir şey yazma!} \\
\\ 
\texttt{Toplamda 8 cezanın 1 kadarını biriktirdiniz.Lütfen tekrar deneyin!} \\
            }
        }
    }
    & \\ \\

    \theutterance \stepcounter{utterance}  
    & \multicolumn{4}{p{0.6\linewidth}}{
        \cellcolor[rgb]{0.9,0.9,0.9}{
            \makecell[{{p{\linewidth}}}]{
                \texttt{\tiny{[P1$\rangle$GM]}}
                \texttt{SÖYLE: Selam! Ortak bir hedef belirlemek için nesne konumlarımızı paylaşalım. Benim konumlarım: C (2,6), L (2,5), P (7,3). Senin konumların neler?} \\
            }
        }
    }
    & & \\ \\

    \theutterance \stepcounter{utterance}  
    & & & \multicolumn{2}{p{0.3\linewidth}}{
        \cellcolor[rgb]{0.9,0.9,0.9}{
            \makecell[{{p{\linewidth}}}]{
                \texttt{\tiny{[GM$|$GM]}}
                \texttt{SÖYLE: Selam! Ortak bir hedef belirlemek için nesne konumlarımızı paylaşalım. Benim konumlarım: C (2,6), L (2,5), P (7,3). Senin konumların neler?} \\
            }
        }
    }
    & & \\ \\

    \theutterance \stepcounter{utterance}  
    & & \multicolumn{4}{p{0.6\linewidth}}{
        \cellcolor[rgb]{0.9,0.9,0.9}{
            \makecell[{{p{\linewidth}}}]{
                \texttt{\tiny{[GM$\rangle$P2]}}
                \texttt{Ben senin oyun yöneticinim ve aşağıdaki ızgarayı oyun tahtası olarak kullandığınız işbirlikçi bir oyun oynuyorsunuz.} \\
\texttt{```} \\
\texttt{ 1234567} \\
\texttt{╔═╤═══╤═╗} \\
\texttt{║◌│◌C◌│◌║ 1} \\
\texttt{╟─┼──┬┘◌║ 2} \\
\texttt{║◌│◌◌│◌◌║ 3} \\
\texttt{╟─┴──┘◌◌║ 4} \\
\texttt{║◌◌◌◌◌◌◌║ 5} \\
\texttt{║◌◌◌◌◌◌◌║ 6} \\
\texttt{║◌◌◌P◌◌L║ 7} \\
\texttt{╚═══════╝} \\
\\ 
\texttt{```} \\
\\ 
\texttt{* Üst kenarda x koordinatları sağa doğru artar, sağ kenarda ise y koordinatları aşağı doğru artar.} \\
\texttt{* Aşağıdaki nesneler oyun tahtasına rastgele yerleştirilmiştir: 'C', 'L', 'P'.} \\
\\ 
\\ 
\texttt{Diğer oyuncu, nesnelerin farklı rastgele konumlara yerleştirildiği oyun tahtasının bir varyasyonunu görür. Sen diğer oyuncunun tahtasını göremezsin ve onlar da seninkini göremez.} \\
\texttt{**Amaç:**} \\
\\ 
\texttt{Her iki oyuncunun da, kendi oyun tahtalarındaki nesneleri hareket ettirerek aynı türden nesnelerin aynı koordinatlara gelmesini sağlaması gerekir. Ortak bir hedef düzen üzerinde anlaşmak için diğer oyuncuyla iletişim kurmalısınız.} \\
\\ 
\\ 
\texttt{**Kurallar:**} \\
\\ 
\texttt{* Her turda, aşağıdaki iki komuttan yalnızca birini gönderebilirsin:} \\
\texttt{ 1. `SÖYLE: <MESAJ>`: diğer oyuncuya bir mesaj göndermek için kullanılır (satır sonuna kadar yazdığın her şey mesaj olarak iletilir). Ben bu mesajı partnerine ileteceğim.} \\
\texttt{2. `HAMLE: <NESNE>, (<X>, <Y>)`: bir nesneyi yeni bir konuma taşımak için kullanılır. Burada <X> sütunu, <Y> ise satırı ifade eder. Yaptığın hamlenin başarılı olup olmadığını sana söyleyeceğim.} \\
\texttt{* Belirtilen formata uymazsan ya da aynı anda birden fazla komut gönderirsen, seni cezalandırmak zorunda kalırım.} \\
\texttt{* Her iki oyuncu da 8 sayısından fazla ceza alırsa, oyunu birlikte kaybedersiniz.} \\
\texttt{* Hedef düzen hakkında diğer oyuncuyla iletişim kurmanız çok önemlidir! Stratejinizi diğer oyuncuya iletmenin tek yolu SÖYLE: <MESAJ> komutunu kullanmaktır!} \\
\\ 
\\ 
\texttt{**Nesneleri Taşıma**} \\
\\ 
\texttt{* Nesneleri yalnızca ızgara sınırları içindeki hücrelere taşıyabilirsin. Hedef hücre boş olmalıdır; yani yalnızca ‘◌’ sembolünü içermelidir.} \\
\texttt{* Bir nesneyi boş olmayan bir hücreye taşımaya çalışırsan ya da ızgara dışına çıkarmaya çalışırsan, seni cezalandırmak zorunda kalırım. Ardından tekrar deneme hakkın olur.} \\
\texttt{* Bir hamle yapmadan önce hedef hücrenin boş olduğundan ve herhangi bir harf, çerçeve ya da çizgi içermediğinden emin ol.} \\
\\ 
\\ 
\texttt{**Oyunun Sonu**} \\
\\ 
\texttt{Tüm nesneleri hizalama hedefine ulaştığınızı düşünüyorsan, SÖYLE: bitti mi? göndererek diğer oyuncudan oyunu bitirmesini isteyebilirsin. Eğer diğer oyuncu senden oyunu bitirmeni isterse ve sen SÖYLE: bitti! ile yanıt verirsen, oyun sona erer.} \\
\\ 
\texttt{Oyun, 12 tur içinde sona ererse her iki oyuncu da kazanır. Bir tur, iki oyuncunun da geçerli bir komut göndermesi olarak tanımlanır.} \\
\\ 
\texttt{**Puanlama:**} \\
\\ 
\texttt{Aynı nesneler her iki oyun tahtasında ne kadar yakın konumlanırsa, o kadar fazla puan kazanırsınız. Cezalar puanınızı düşürür. Rekoru kırabilir misiniz?} \\
\\ 
\texttt{Diğer oyuncu şu mesajı göndererek oyunu başlattı:} \\
\\ 
\texttt{"Selam! Ortak bir hedef belirlemek için nesne konumlarımızı paylaşalım. Benim konumlarım: C (2,6), L (2,5), P (7,3). Senin konumların neler?"} \\
\\ 
\texttt{İlk komutun ne?} \\
            }
        }
    }
    & \\ \\

    \theutterance \stepcounter{utterance}  
    & & & \multicolumn{4}{p{0.6\linewidth}}{
        \cellcolor[rgb]{0.9,0.9,0.9}{
            \makecell[{{p{\linewidth}}}]{
                \texttt{\tiny{[GM$\langle$P2]}}
                \texttt{SÖYLE: Selam! Benim konumlarım: C (4,1), P (4,7), L (7,7). Hedef olarak tüm nesneleri (C, P, L) ızgaranın orta kısmında veya belirli bir köşede toplamaya ne dersin? Örneğin, C (4,3), P (4,4), L (4,5) gibi bir dikey sıra belirleyebiliriz. Senin önerin nedir?} \\
            }
        }
    }
     \\ \\

    \theutterance \stepcounter{utterance}  
    & & & \multicolumn{2}{p{0.3\linewidth}}{
        \cellcolor[rgb]{0.9,0.9,0.9}{
            \makecell[{{p{\linewidth}}}]{
                \texttt{\tiny{[GM$|$GM]}}
                \texttt{SÖYLE: Selam! Benim konumlarım: C (4,1), P (4,7), L (7,7). Hedef olarak tüm nesneleri (C, P, L) ızgaranın orta kısmında veya belirli bir köşede toplamaya ne dersin? Örneğin, C (4,3), P (4,4), L (4,5) gibi bir dikey sıra belirleyebiliriz. Senin önerin nedir?} \\
            }
        }
    }
    & & \\ \\

    \theutterance \stepcounter{utterance}  
    & & \multicolumn{4}{p{0.6\linewidth}}{
        \cellcolor[rgb]{0.9,0.9,0.9}{
            \makecell[{{p{\linewidth}}}]{
                \texttt{\tiny{[P1$\langle$GM]}}
                \texttt{Mesajın diğer oyuncuya iletildi.} \\
\\ 
\\ 
\texttt{Şu anda en fazla 12 turdan 1. turu oynuyorsunuz.} \\
\texttt{Toplamda 8 cezanın 1 kadarını biriktirdiniz.Diğer oyuncu şu mesajı gönderdi:} \\
\\ 
\\ 
\texttt{"Selam! Benim konumlarım: C (4,1), P (4,7), L (7,7). Hedef olarak tüm nesneleri (C, P, L) ızgaranın orta kısmında veya belirli bir köşede toplamaya ne dersin? Örneğin, C (4,3), P (4,4), L (4,5) gibi bir dikey sıra belirleyebiliriz. Senin önerin nedir?"} \\
\\ 
\\ 
\texttt{Bir sonraki komutun nedir?} \\
            }
        }
    }
    & \\ \\

    \theutterance \stepcounter{utterance}  
    & \multicolumn{4}{p{0.6\linewidth}}{
        \cellcolor[rgb]{0.9,0.9,0.9}{
            \makecell[{{p{\linewidth}}}]{
                \texttt{\tiny{[P1$\rangle$GM]}}
                \texttt{SÖYLE: Önerin harika! C (4,3), P (4,4) ve L (4,5) koordinatları çok mantıklı bir hedef. Ben de bu hedefi kabul ediyorum. Ben hareket etmeye başlıyorum, sen de kendi tahtandaki nesneleri bu koordinatlara taşımaya başla. İlk hamlem olarak C'yi taşıyorum.} \\
            }
        }
    }
    & & \\ \\

    \theutterance \stepcounter{utterance}  
    & & & \multicolumn{2}{p{0.3\linewidth}}{
        \cellcolor[rgb]{0.9,0.9,0.9}{
            \makecell[{{p{\linewidth}}}]{
                \texttt{\tiny{[GM$|$GM]}}
                \texttt{SÖYLE: Önerin harika! C (4,3), P (4,4) ve L (4,5) koordinatları çok mantıklı bir hedef. Ben de bu hedefi kabul ediyorum. Ben hareket etmeye başlıyorum, sen de kendi tahtandaki nesneleri bu koordinatlara taşımaya başla. İlk hamlem olarak C'yi taşıyorum.} \\
            }
        }
    }
    & & \\ \\

    \theutterance \stepcounter{utterance}  
    & & \multicolumn{4}{p{0.6\linewidth}}{
        \cellcolor[rgb]{0.9,0.9,0.9}{
            \makecell[{{p{\linewidth}}}]{
                \texttt{\tiny{[GM$\rangle$P2]}}
                \texttt{Mesajın diğer oyuncuya iletildi.} \\
\\ 
\\ 
\texttt{Şu anda en fazla 12 turdan 2. turu oynuyorsunuz.} \\
\texttt{Toplamda 8 cezanın 1 kadarını biriktirdiniz.Diğer oyuncu şu mesajı gönderdi:} \\
\\ 
\\ 
\texttt{"Önerin harika! C (4,3), P (4,4) ve L (4,5) koordinatları çok mantıklı bir hedef. Ben de bu hedefi kabul ediyorum. Ben hareket etmeye başlıyorum, sen de kendi tahtandaki nesneleri bu koordinatlara taşımaya başla. İlk hamlem olarak C'yi taşıyorum."} \\
\\ 
\\ 
\texttt{Bir sonraki komutun nedir?} \\
            }
        }
    }
    & \\ \\

    \theutterance \stepcounter{utterance}  
    & & & \multicolumn{4}{p{0.6\linewidth}}{
        \cellcolor[rgb]{0.9,0.9,0.9}{
            \makecell[{{p{\linewidth}}}]{
                \texttt{\tiny{[GM$\langle$P2]}}
                \texttt{HAMLE: C, (4,3)} \\
            }
        }
    }
     \\ \\

    \theutterance \stepcounter{utterance}  
    & & & \multicolumn{2}{p{0.3\linewidth}}{
        \cellcolor[rgb]{0.9,0.9,0.9}{
            \makecell[{{p{\linewidth}}}]{
                \texttt{\tiny{[GM$|$GM]}}
                \texttt{HAMLE: C, (4,3)} \\
            }
        }
    }
    & & \\ \\

    \theutterance \stepcounter{utterance}  
    & & & \multicolumn{2}{p{0.3\linewidth}}{
        \cellcolor[rgb]{0.9,0.9,0.9}{
            \makecell[{{p{\linewidth}}}]{
                \texttt{\tiny{[GM$|$GM]}}
                \texttt{C moved to (4, 3)} \\
            }
        }
    }
    & & \\ \\

    \theutterance \stepcounter{utterance}  
    & & \multicolumn{4}{p{0.6\linewidth}}{
        \cellcolor[rgb]{0.9,0.9,0.9}{
            \makecell[{{p{\linewidth}}}]{
                \texttt{\tiny{[P1$\langle$GM]}}
                \texttt{Mesajın diğer oyuncuya iletildi.} \\
\\ 
\\ 
\texttt{Şu anda en fazla 12 turdan 2. turu oynuyorsunuz.} \\
\texttt{Toplamda 8 cezanın 1 kadarını biriktirdiniz.Diğer oyuncu kendi ızgarasında bir nesneyi hareket ettirdi.} \\
\\ 
\\ 
\texttt{Bir sonraki komutun nedir?} \\
            }
        }
    }
    & \\ \\

    \theutterance \stepcounter{utterance}  
    & \multicolumn{4}{p{0.6\linewidth}}{
        \cellcolor[rgb]{0.9,0.9,0.9}{
            \makecell[{{p{\linewidth}}}]{
                \texttt{\tiny{[P1$\rangle$GM]}}
                \texttt{HAMLE: C, (4, 3)} \\
            }
        }
    }
    & & \\ \\

    \theutterance \stepcounter{utterance}  
    & & & \multicolumn{2}{p{0.3\linewidth}}{
        \cellcolor[rgb]{0.9,0.9,0.9}{
            \makecell[{{p{\linewidth}}}]{
                \texttt{\tiny{[GM$|$GM]}}
                \texttt{HAMLE: C, (4, 3)} \\
            }
        }
    }
    & & \\ \\

    \theutterance \stepcounter{utterance}  
    & & & \multicolumn{2}{p{0.3\linewidth}}{
        \cellcolor[rgb]{0.9,0.9,0.9}{
            \makecell[{{p{\linewidth}}}]{
                \texttt{\tiny{[GM$|$GM]}}
                \texttt{C moved to (4, 3)} \\
            }
        }
    }
    & & \\ \\

    \theutterance \stepcounter{utterance}  
    & & \multicolumn{4}{p{0.6\linewidth}}{
        \cellcolor[rgb]{0.9,0.9,0.9}{
            \makecell[{{p{\linewidth}}}]{
                \texttt{\tiny{[GM$\rangle$P2]}}
                \texttt{C’ başarıyla (4,3) koordinatlarına taşındı. Güncellenmiş ızgaran şu şekilde:} \\
\texttt{```} \\
\texttt{ 1234567} \\
\texttt{╔═╤═══╤═╗} \\
\texttt{║◌│◌◌◌│◌║ 1} \\
\texttt{╟─┼──┬┘◌║ 2} \\
\texttt{║◌│◌C│◌◌║ 3} \\
\texttt{╟─┴──┘◌◌║ 4} \\
\texttt{║◌◌◌◌◌◌◌║ 5} \\
\texttt{║◌◌◌◌◌◌◌║ 6} \\
\texttt{║◌◌◌P◌◌L║ 7} \\
\texttt{╚═══════╝} \\
\texttt{```} \\
\texttt{Şu anda en fazla 12 turdan 3. turu oynuyorsunuz.} \\
\texttt{Toplamda 8 cezanın 1 kadarını biriktirdiniz.Diğer oyuncu kendi ızgarasında bir nesneyi hareket ettirdi.} \\
\\ 
\\ 
\texttt{Bir sonraki komutun nedir?} \\
            }
        }
    }
    & \\ \\

    \theutterance \stepcounter{utterance}  
    & & & \multicolumn{4}{p{0.6\linewidth}}{
        \cellcolor[rgb]{0.9,0.9,0.9}{
            \makecell[{{p{\linewidth}}}]{
                \texttt{\tiny{[GM$\langle$P2]}}
                \texttt{HAMLE: P, (4,4)} \\
            }
        }
    }
     \\ \\

    \theutterance \stepcounter{utterance}  
    & & & \multicolumn{2}{p{0.3\linewidth}}{
        \cellcolor[rgb]{0.9,0.9,0.9}{
            \makecell[{{p{\linewidth}}}]{
                \texttt{\tiny{[GM$|$GM]}}
                \texttt{HAMLE: P, (4,4)} \\
            }
        }
    }
    & & \\ \\

    \theutterance \stepcounter{utterance}  
    & & & \multicolumn{2}{p{0.3\linewidth}}{
        \cellcolor[rgb]{0.9,0.9,0.9}{
            \makecell[{{p{\linewidth}}}]{
                \texttt{\tiny{[GM$|$GM]}}
                \texttt{Ceza: (4,4) boş değil; "─" içeriyor.} \\
\\ 
\texttt{Toplamda 8 cezanın 2 kadarını biriktirdiniz.Lütfen tekrar deneyin!} \\
            }
        }
    }
    & & \\ \\

    \theutterance \stepcounter{utterance}  
    & & \multicolumn{4}{p{0.6\linewidth}}{
        \cellcolor[rgb]{0.9,0.9,0.9}{
            \makecell[{{p{\linewidth}}}]{
                \texttt{\tiny{[GM$\rangle$P2]}}
                \texttt{Ceza: (4,4) boş değil; "─" içeriyor.} \\
\\ 
\texttt{Toplamda 8 cezanın 2 kadarını biriktirdiniz.Lütfen tekrar deneyin!} \\
            }
        }
    }
    & \\ \\

    \theutterance \stepcounter{utterance}  
    & & & \multicolumn{4}{p{0.6\linewidth}}{
        \cellcolor[rgb]{0.9,0.9,0.9}{
            \makecell[{{p{\linewidth}}}]{
                \texttt{\tiny{[GM$\langle$P2]}}
                \texttt{HAMLE: P, (4,5)} \\
            }
        }
    }
     \\ \\

    \theutterance \stepcounter{utterance}  
    & & & \multicolumn{2}{p{0.3\linewidth}}{
        \cellcolor[rgb]{0.9,0.9,0.9}{
            \makecell[{{p{\linewidth}}}]{
                \texttt{\tiny{[GM$|$GM]}}
                \texttt{HAMLE: P, (4,5)} \\
            }
        }
    }
    & & \\ \\

    \theutterance \stepcounter{utterance}  
    & & & \multicolumn{2}{p{0.3\linewidth}}{
        \cellcolor[rgb]{0.9,0.9,0.9}{
            \makecell[{{p{\linewidth}}}]{
                \texttt{\tiny{[GM$|$GM]}}
                \texttt{P moved to (4, 5)} \\
            }
        }
    }
    & & \\ \\

    \theutterance \stepcounter{utterance}  
    & & \multicolumn{4}{p{0.6\linewidth}}{
        \cellcolor[rgb]{0.9,0.9,0.9}{
            \makecell[{{p{\linewidth}}}]{
                \texttt{\tiny{[P1$\langle$GM]}}
                \texttt{C’ başarıyla (4,3) koordinatlarına taşındı. Güncellenmiş ızgaran şu şekilde:} \\
\texttt{```} \\
\texttt{ 1234567} \\
\texttt{╔═╤═══╤═╗} \\
\texttt{║◌│◌◌◌│◌║ 1} \\
\texttt{╟─┼──┬┘◌║ 2} \\
\texttt{║◌│◌C│◌P║ 3} \\
\texttt{╟─┴──┘◌◌║ 4} \\
\texttt{║◌L◌◌◌◌◌║ 5} \\
\texttt{║◌◌◌◌◌◌◌║ 6} \\
\texttt{║◌◌◌◌◌◌◌║ 7} \\
\texttt{╚═══════╝} \\
\texttt{```} \\
\texttt{Şu anda en fazla 12 turdan 3. turu oynuyorsunuz.} \\
\texttt{Toplamda 8 cezanın 2 kadarını biriktirdiniz.Diğer oyuncu kendi ızgarasında bir nesneyi hareket ettirdi.} \\
\\ 
\\ 
\texttt{Bir sonraki komutun nedir?} \\
            }
        }
    }
    & \\ \\

    \theutterance \stepcounter{utterance}  
    & \multicolumn{4}{p{0.6\linewidth}}{
        \cellcolor[rgb]{0.9,0.9,0.9}{
            \makecell[{{p{\linewidth}}}]{
                \texttt{\tiny{[P1$\rangle$GM]}}
                \texttt{HAMLE: L, (4, 5)} \\
            }
        }
    }
    & & \\ \\

    \theutterance \stepcounter{utterance}  
    & & & \multicolumn{2}{p{0.3\linewidth}}{
        \cellcolor[rgb]{0.9,0.9,0.9}{
            \makecell[{{p{\linewidth}}}]{
                \texttt{\tiny{[GM$|$GM]}}
                \texttt{HAMLE: L, (4, 5)} \\
            }
        }
    }
    & & \\ \\

    \theutterance \stepcounter{utterance}  
    & & & \multicolumn{2}{p{0.3\linewidth}}{
        \cellcolor[rgb]{0.9,0.9,0.9}{
            \makecell[{{p{\linewidth}}}]{
                \texttt{\tiny{[GM$|$GM]}}
                \texttt{L moved to (4, 5)} \\
            }
        }
    }
    & & \\ \\

    \theutterance \stepcounter{utterance}  
    & & \multicolumn{4}{p{0.6\linewidth}}{
        \cellcolor[rgb]{0.9,0.9,0.9}{
            \makecell[{{p{\linewidth}}}]{
                \texttt{\tiny{[GM$\rangle$P2]}}
                \texttt{P’ başarıyla (4,5) koordinatlarına taşındı. Güncellenmiş ızgaran şu şekilde:} \\
\texttt{```} \\
\texttt{ 1234567} \\
\texttt{╔═╤═══╤═╗} \\
\texttt{║◌│◌◌◌│◌║ 1} \\
\texttt{╟─┼──┬┘◌║ 2} \\
\texttt{║◌│◌C│◌◌║ 3} \\
\texttt{╟─┴──┘◌◌║ 4} \\
\texttt{║◌◌◌P◌◌◌║ 5} \\
\texttt{║◌◌◌◌◌◌◌║ 6} \\
\texttt{║◌◌◌◌◌◌L║ 7} \\
\texttt{╚═══════╝} \\
\texttt{```} \\
\texttt{Şu anda en fazla 12 turdan 4. turu oynuyorsunuz.} \\
\texttt{Toplamda 8 cezanın 2 kadarını biriktirdiniz.Diğer oyuncu kendi ızgarasında bir nesneyi hareket ettirdi.} \\
\\ 
\\ 
\texttt{Bir sonraki komutun nedir?} \\
            }
        }
    }
    & \\ \\

    \theutterance \stepcounter{utterance}  
    & & & \multicolumn{4}{p{0.6\linewidth}}{
        \cellcolor[rgb]{0.9,0.9,0.9}{
            \makecell[{{p{\linewidth}}}]{
                \texttt{\tiny{[GM$\langle$P2]}}
                \texttt{HAMLE: L, (4,5)} \\
            }
        }
    }
     \\ \\

    \theutterance \stepcounter{utterance}  
    & & & \multicolumn{2}{p{0.3\linewidth}}{
        \cellcolor[rgb]{0.9,0.9,0.9}{
            \makecell[{{p{\linewidth}}}]{
                \texttt{\tiny{[GM$|$GM]}}
                \texttt{HAMLE: L, (4,5)} \\
            }
        }
    }
    & & \\ \\

    \theutterance \stepcounter{utterance}  
    & & & \multicolumn{2}{p{0.3\linewidth}}{
        \cellcolor[rgb]{0.9,0.9,0.9}{
            \makecell[{{p{\linewidth}}}]{
                \texttt{\tiny{[GM$|$GM]}}
                \texttt{Ceza: (4,5) boş değil; "P" içeriyor.} \\
\\ 
\texttt{Toplamda 8 cezanın 3 kadarını biriktirdiniz.Lütfen tekrar deneyin!} \\
            }
        }
    }
    & & \\ \\

    \theutterance \stepcounter{utterance}  
    & & \multicolumn{4}{p{0.6\linewidth}}{
        \cellcolor[rgb]{0.9,0.9,0.9}{
            \makecell[{{p{\linewidth}}}]{
                \texttt{\tiny{[GM$\rangle$P2]}}
                \texttt{Ceza: (4,5) boş değil; "P" içeriyor.} \\
\\ 
\texttt{Toplamda 8 cezanın 3 kadarını biriktirdiniz.Lütfen tekrar deneyin!} \\
            }
        }
    }
    & \\ \\

    \theutterance \stepcounter{utterance}  
    & & & \multicolumn{4}{p{0.6\linewidth}}{
        \cellcolor[rgb]{0.9,0.9,0.9}{
            \makecell[{{p{\linewidth}}}]{
                \texttt{\tiny{[GM$\langle$P2]}}
                \texttt{HAMLE: L, (4,6)} \\
            }
        }
    }
     \\ \\

    \theutterance \stepcounter{utterance}  
    & & & \multicolumn{2}{p{0.3\linewidth}}{
        \cellcolor[rgb]{0.9,0.9,0.9}{
            \makecell[{{p{\linewidth}}}]{
                \texttt{\tiny{[GM$|$GM]}}
                \texttt{HAMLE: L, (4,6)} \\
            }
        }
    }
    & & \\ \\

    \theutterance \stepcounter{utterance}  
    & & & \multicolumn{2}{p{0.3\linewidth}}{
        \cellcolor[rgb]{0.9,0.9,0.9}{
            \makecell[{{p{\linewidth}}}]{
                \texttt{\tiny{[GM$|$GM]}}
                \texttt{L moved to (4, 6)} \\
            }
        }
    }
    & & \\ \\

    \theutterance \stepcounter{utterance}  
    & & \multicolumn{4}{p{0.6\linewidth}}{
        \cellcolor[rgb]{0.9,0.9,0.9}{
            \makecell[{{p{\linewidth}}}]{
                \texttt{\tiny{[P1$\langle$GM]}}
                \texttt{L’ başarıyla (4,5) koordinatlarına taşındı. Güncellenmiş ızgaran şu şekilde:} \\
\texttt{```} \\
\texttt{ 1234567} \\
\texttt{╔═╤═══╤═╗} \\
\texttt{║◌│◌◌◌│◌║ 1} \\
\texttt{╟─┼──┬┘◌║ 2} \\
\texttt{║◌│◌C│◌P║ 3} \\
\texttt{╟─┴──┘◌◌║ 4} \\
\texttt{║◌◌◌L◌◌◌║ 5} \\
\texttt{║◌◌◌◌◌◌◌║ 6} \\
\texttt{║◌◌◌◌◌◌◌║ 7} \\
\texttt{╚═══════╝} \\
\texttt{```} \\
\texttt{Şu anda en fazla 12 turdan 4. turu oynuyorsunuz.} \\
\texttt{Toplamda 8 cezanın 3 kadarını biriktirdiniz.Diğer oyuncu kendi ızgarasında bir nesneyi hareket ettirdi.} \\
\\ 
\\ 
\texttt{Bir sonraki komutun nedir?} \\
            }
        }
    }
    & \\ \\

    \theutterance \stepcounter{utterance}  
    & \multicolumn{4}{p{0.6\linewidth}}{
        \cellcolor[rgb]{0.9,0.9,0.9}{
            \makecell[{{p{\linewidth}}}]{
                \texttt{\tiny{[P1$\rangle$GM]}}
                \texttt{HAMLE: P, (4, 4)} \\
            }
        }
    }
    & & \\ \\

    \theutterance \stepcounter{utterance}  
    & & & \multicolumn{2}{p{0.3\linewidth}}{
        \cellcolor[rgb]{0.9,0.9,0.9}{
            \makecell[{{p{\linewidth}}}]{
                \texttt{\tiny{[GM$|$GM]}}
                \texttt{HAMLE: P, (4, 4)} \\
            }
        }
    }
    & & \\ \\

    \theutterance \stepcounter{utterance}  
    & & & \multicolumn{2}{p{0.3\linewidth}}{
        \cellcolor[rgb]{0.9,0.9,0.9}{
            \makecell[{{p{\linewidth}}}]{
                \texttt{\tiny{[GM$|$GM]}}
                \texttt{Ceza: (4,4) boş değil; "─" içeriyor.} \\
\\ 
\texttt{Toplamda 8 cezanın 4 kadarını biriktirdiniz.Lütfen tekrar deneyin!} \\
            }
        }
    }
    & & \\ \\

    \theutterance \stepcounter{utterance}  
    & & \multicolumn{4}{p{0.6\linewidth}}{
        \cellcolor[rgb]{0.9,0.9,0.9}{
            \makecell[{{p{\linewidth}}}]{
                \texttt{\tiny{[P1$\langle$GM]}}
                \texttt{Ceza: (4,4) boş değil; "─" içeriyor.} \\
\\ 
\texttt{Toplamda 8 cezanın 4 kadarını biriktirdiniz.Lütfen tekrar deneyin!} \\
            }
        }
    }
    & \\ \\

    \theutterance \stepcounter{utterance}  
    & \multicolumn{4}{p{0.6\linewidth}}{
        \cellcolor[rgb]{0.9,0.9,0.9}{
            \makecell[{{p{\linewidth}}}]{
                \texttt{\tiny{[P1$\rangle$GM]}}
                \texttt{HAMLE: P, (7, 4)} \\
            }
        }
    }
    & & \\ \\

    \theutterance \stepcounter{utterance}  
    & & & \multicolumn{2}{p{0.3\linewidth}}{
        \cellcolor[rgb]{0.9,0.9,0.9}{
            \makecell[{{p{\linewidth}}}]{
                \texttt{\tiny{[GM$|$GM]}}
                \texttt{HAMLE: P, (7, 4)} \\
            }
        }
    }
    & & \\ \\

    \theutterance \stepcounter{utterance}  
    & & & \multicolumn{2}{p{0.3\linewidth}}{
        \cellcolor[rgb]{0.9,0.9,0.9}{
            \makecell[{{p{\linewidth}}}]{
                \texttt{\tiny{[GM$|$GM]}}
                \texttt{P moved to (7, 4)} \\
            }
        }
    }
    & & \\ \\

    \theutterance \stepcounter{utterance}  
    & & \multicolumn{4}{p{0.6\linewidth}}{
        \cellcolor[rgb]{0.9,0.9,0.9}{
            \makecell[{{p{\linewidth}}}]{
                \texttt{\tiny{[GM$\rangle$P2]}}
                \texttt{L’ başarıyla (4,6) koordinatlarına taşındı. Güncellenmiş ızgaran şu şekilde:} \\
\texttt{```} \\
\texttt{ 1234567} \\
\texttt{╔═╤═══╤═╗} \\
\texttt{║◌│◌◌◌│◌║ 1} \\
\texttt{╟─┼──┬┘◌║ 2} \\
\texttt{║◌│◌C│◌◌║ 3} \\
\texttt{╟─┴──┘◌◌║ 4} \\
\texttt{║◌◌◌P◌◌◌║ 5} \\
\texttt{║◌◌◌L◌◌◌║ 6} \\
\texttt{║◌◌◌◌◌◌◌║ 7} \\
\texttt{╚═══════╝} \\
\texttt{```} \\
\texttt{Şu anda en fazla 12 turdan 5. turu oynuyorsunuz.} \\
\texttt{Toplamda 8 cezanın 4 kadarını biriktirdiniz.Diğer oyuncu kendi ızgarasında bir nesneyi hareket ettirdi.} \\
\\ 
\\ 
\texttt{Bir sonraki komutun nedir?} \\
            }
        }
    }
    & \\ \\

    \theutterance \stepcounter{utterance}  
    & & & \multicolumn{4}{p{0.6\linewidth}}{
        \cellcolor[rgb]{0.9,0.9,0.9}{
            \makecell[{{p{\linewidth}}}]{
                \texttt{\tiny{[GM$\langle$P2]}}
                \texttt{HAMLE: L, (4,5)} \\
            }
        }
    }
     \\ \\

    \theutterance \stepcounter{utterance}  
    & & & \multicolumn{2}{p{0.3\linewidth}}{
        \cellcolor[rgb]{0.9,0.9,0.9}{
            \makecell[{{p{\linewidth}}}]{
                \texttt{\tiny{[GM$|$GM]}}
                \texttt{HAMLE: L, (4,5)} \\
            }
        }
    }
    & & \\ \\

    \theutterance \stepcounter{utterance}  
    & & & \multicolumn{2}{p{0.3\linewidth}}{
        \cellcolor[rgb]{0.9,0.9,0.9}{
            \makecell[{{p{\linewidth}}}]{
                \texttt{\tiny{[GM$|$GM]}}
                \texttt{Ceza: (4,5) boş değil; "P" içeriyor.} \\
\\ 
\texttt{Toplamda 8 cezanın 5 kadarını biriktirdiniz.Lütfen tekrar deneyin!} \\
            }
        }
    }
    & & \\ \\

    \theutterance \stepcounter{utterance}  
    & & \multicolumn{4}{p{0.6\linewidth}}{
        \cellcolor[rgb]{0.9,0.9,0.9}{
            \makecell[{{p{\linewidth}}}]{
                \texttt{\tiny{[GM$\rangle$P2]}}
                \texttt{Ceza: (4,5) boş değil; "P" içeriyor.} \\
\\ 
\texttt{Toplamda 8 cezanın 5 kadarını biriktirdiniz.Lütfen tekrar deneyin!} \\
            }
        }
    }
    & \\ \\

    \theutterance \stepcounter{utterance}  
    & & & \multicolumn{4}{p{0.6\linewidth}}{
        \cellcolor[rgb]{0.9,0.9,0.9}{
            \makecell[{{p{\linewidth}}}]{
                \texttt{\tiny{[GM$\langle$P2]}}
                \texttt{HAMLE: L, (4,7)} \\
            }
        }
    }
     \\ \\

    \theutterance \stepcounter{utterance}  
    & & & \multicolumn{2}{p{0.3\linewidth}}{
        \cellcolor[rgb]{0.9,0.9,0.9}{
            \makecell[{{p{\linewidth}}}]{
                \texttt{\tiny{[GM$|$GM]}}
                \texttt{HAMLE: L, (4,7)} \\
            }
        }
    }
    & & \\ \\

    \theutterance \stepcounter{utterance}  
    & & & \multicolumn{2}{p{0.3\linewidth}}{
        \cellcolor[rgb]{0.9,0.9,0.9}{
            \makecell[{{p{\linewidth}}}]{
                \texttt{\tiny{[GM$|$GM]}}
                \texttt{L moved to (4, 7)} \\
            }
        }
    }
    & & \\ \\

    \theutterance \stepcounter{utterance}  
    & & \multicolumn{4}{p{0.6\linewidth}}{
        \cellcolor[rgb]{0.9,0.9,0.9}{
            \makecell[{{p{\linewidth}}}]{
                \texttt{\tiny{[P1$\langle$GM]}}
                \texttt{P’ başarıyla (7,4) koordinatlarına taşındı. Güncellenmiş ızgaran şu şekilde:} \\
\texttt{```} \\
\texttt{ 1234567} \\
\texttt{╔═╤═══╤═╗} \\
\texttt{║◌│◌◌◌│◌║ 1} \\
\texttt{╟─┼──┬┘◌║ 2} \\
\texttt{║◌│◌C│◌◌║ 3} \\
\texttt{╟─┴──┘◌P║ 4} \\
\texttt{║◌◌◌L◌◌◌║ 5} \\
\texttt{║◌◌◌◌◌◌◌║ 6} \\
\texttt{║◌◌◌◌◌◌◌║ 7} \\
\texttt{╚═══════╝} \\
\texttt{```} \\
\texttt{Şu anda en fazla 12 turdan 5. turu oynuyorsunuz.} \\
\texttt{Toplamda 8 cezanın 5 kadarını biriktirdiniz.Diğer oyuncu kendi ızgarasında bir nesneyi hareket ettirdi.} \\
\\ 
\\ 
\texttt{Bir sonraki komutun nedir?} \\
            }
        }
    }
    & \\ \\

    \theutterance \stepcounter{utterance}  
    & \multicolumn{4}{p{0.6\linewidth}}{
        \cellcolor[rgb]{0.9,0.9,0.9}{
            \makecell[{{p{\linewidth}}}]{
                \texttt{\tiny{[P1$\rangle$GM]}}
                \texttt{SÖYLE: P'yi (4,4) koordinatına taşımaya çalışırken hata yaptım, özür dilerim. Hedefimiz olan (4,4) hücresi aslında bir çizgi içeriyor gibi görünüyor, bu yüzden P'yi (4,4) yerine (5,4) veya (6,4) gibi boş bir yere çekip sonra hedefi tekrar değerlendirelim. Ancak şu anki konumumda P'yi (4,4) yerine (5,4) koordinatına taşımayı deneyeceğim.} \\
            }
        }
    }
    & & \\ \\

    \theutterance \stepcounter{utterance}  
    & & & \multicolumn{2}{p{0.3\linewidth}}{
        \cellcolor[rgb]{0.9,0.9,0.9}{
            \makecell[{{p{\linewidth}}}]{
                \texttt{\tiny{[GM$|$GM]}}
                \texttt{SÖYLE: P'yi (4,4) koordinatına taşımaya çalışırken hata yaptım, özür dilerim. Hedefimiz olan (4,4) hücresi aslında bir çizgi içeriyor gibi görünüyor, bu yüzden P'yi (4,4) yerine (5,4) veya (6,4) gibi boş bir yere çekip sonra hedefi tekrar değerlendirelim. Ancak şu anki konumumda P'yi (4,4) yerine (5,4) koordinatına taşımayı deneyeceğim.} \\
            }
        }
    }
    & & \\ \\

    \theutterance \stepcounter{utterance}  
    & & \multicolumn{4}{p{0.6\linewidth}}{
        \cellcolor[rgb]{0.9,0.9,0.9}{
            \makecell[{{p{\linewidth}}}]{
                \texttt{\tiny{[GM$\rangle$P2]}}
                \texttt{L’ başarıyla (4,7) koordinatlarına taşındı. Güncellenmiş ızgaran şu şekilde:} \\
\texttt{```} \\
\texttt{ 1234567} \\
\texttt{╔═╤═══╤═╗} \\
\texttt{║◌│◌◌◌│◌║ 1} \\
\texttt{╟─┼──┬┘◌║ 2} \\
\texttt{║◌│◌C│◌◌║ 3} \\
\texttt{╟─┴──┘◌◌║ 4} \\
\texttt{║◌◌◌P◌◌◌║ 5} \\
\texttt{║◌◌◌◌◌◌◌║ 6} \\
\texttt{║◌◌◌L◌◌◌║ 7} \\
\texttt{╚═══════╝} \\
\texttt{```} \\
\texttt{Şu anda en fazla 12 turdan 6. turu oynuyorsunuz.} \\
\texttt{Toplamda 8 cezanın 5 kadarını biriktirdiniz.Diğer oyuncu şu mesajı gönderdi:} \\
\\ 
\\ 
\texttt{"P'yi (4,4) koordinatına taşımaya çalışırken hata yaptım, özür dilerim. Hedefimiz olan (4,4) hücresi aslında bir çizgi içeriyor gibi görünüyor, bu yüzden P'yi (4,4) yerine (5,4) veya (6,4) gibi boş bir yere çekip sonra hedefi tekrar değerlendirelim. Ancak şu anki konumumda P'yi (4,4) yerine (5,4) koordinatına taşımayı deneyeceğim."} \\
\\ 
\\ 
\texttt{Bir sonraki komutun nedir?} \\
            }
        }
    }
    & \\ \\

    \theutterance \stepcounter{utterance}  
    & & & \multicolumn{4}{p{0.6\linewidth}}{
        \cellcolor[rgb]{0.9,0.9,0.9}{
            \makecell[{{p{\linewidth}}}]{
                \texttt{\tiny{[GM$\langle$P2]}}
                \texttt{SÖYLE: Sorun değil, hatalar olabilir. Benim tahtamda (4,4) hücresi de çizgi içeriyor gibi görünüyor, bu yüzden hedefimizi (4,3), (4,5) ve (4,6) gibi daha güvenli, boş hücreler olan bir dikey sıra olarak güncelleyebiliriz. Ben şu an L'yi (4,7) konumuna çektim. P'yi (4,5) konumuna taşımaya ne dersin?} \\
            }
        }
    }
     \\ \\

    \theutterance \stepcounter{utterance}  
    & & & \multicolumn{2}{p{0.3\linewidth}}{
        \cellcolor[rgb]{0.9,0.9,0.9}{
            \makecell[{{p{\linewidth}}}]{
                \texttt{\tiny{[GM$|$GM]}}
                \texttt{SÖYLE: Sorun değil, hatalar olabilir. Benim tahtamda (4,4) hücresi de çizgi içeriyor gibi görünüyor, bu yüzden hedefimizi (4,3), (4,5) ve (4,6) gibi daha güvenli, boş hücreler olan bir dikey sıra olarak güncelleyebiliriz. Ben şu an L'yi (4,7) konumuna çektim. P'yi (4,5) konumuna taşımaya ne dersin?} \\
            }
        }
    }
    & & \\ \\

    \theutterance \stepcounter{utterance}  
    & & \multicolumn{4}{p{0.6\linewidth}}{
        \cellcolor[rgb]{0.9,0.9,0.9}{
            \makecell[{{p{\linewidth}}}]{
                \texttt{\tiny{[P1$\langle$GM]}}
                \texttt{Mesajın diğer oyuncuya iletildi.} \\
\\ 
\\ 
\texttt{Şu anda en fazla 12 turdan 6. turu oynuyorsunuz.} \\
\texttt{Toplamda 8 cezanın 5 kadarını biriktirdiniz.Diğer oyuncu şu mesajı gönderdi:} \\
\\ 
\\ 
\texttt{"Sorun değil, hatalar olabilir. Benim tahtamda (4,4) hücresi de çizgi içeriyor gibi görünüyor, bu yüzden hedefimizi (4,3), (4,5) ve (4,6) gibi daha güvenli, boş hücreler olan bir dikey sıra olarak güncelleyebiliriz. Ben şu an L'yi (4,7) konumuna çektim. P'yi (4,5) konumuna taşımaya ne dersin?"} \\
\\ 
\\ 
\texttt{Bir sonraki komutun nedir?} \\
            }
        }
    }
    & \\ \\

    \theutterance \stepcounter{utterance}  
    & \multicolumn{4}{p{0.6\linewidth}}{
        \cellcolor[rgb]{0.9,0.9,0.9}{
            \makecell[{{p{\linewidth}}}]{
                \texttt{\tiny{[P1$\rangle$GM]}}
                \texttt{SÖYLE: Yeni hedefimizi kabul ediyorum: C (4,3), P (4,5) ve L (4,6). Sen L'yi (4,7)'ye çektiysen, ben de P'yi (4,5) konumuna taşımaya çalışacağım. Ancak benim tahtamda (4,5) hücresinde şu an L var, bu yüzden önce L'yi başka bir yere çekmem gerekecek. Önce L'yi (4,6) konumuna taşımayı deniyorum.} \\
            }
        }
    }
    & & \\ \\

    \theutterance \stepcounter{utterance}  
    & & & \multicolumn{2}{p{0.3\linewidth}}{
        \cellcolor[rgb]{0.9,0.9,0.9}{
            \makecell[{{p{\linewidth}}}]{
                \texttt{\tiny{[GM$|$GM]}}
                \texttt{SÖYLE: Yeni hedefimizi kabul ediyorum: C (4,3), P (4,5) ve L (4,6). Sen L'yi (4,7)'ye çektiysen, ben de P'yi (4,5) konumuna taşımaya çalışacağım. Ancak benim tahtamda (4,5) hücresinde şu an L var, bu yüzden önce L'yi başka bir yere çekmem gerekecek. Önce L'yi (4,6) konumuna taşımayı deniyorum.} \\
            }
        }
    }
    & & \\ \\

    \theutterance \stepcounter{utterance}  
    & & \multicolumn{4}{p{0.6\linewidth}}{
        \cellcolor[rgb]{0.9,0.9,0.9}{
            \makecell[{{p{\linewidth}}}]{
                \texttt{\tiny{[GM$\rangle$P2]}}
                \texttt{Mesajın diğer oyuncuya iletildi.} \\
\\ 
\\ 
\texttt{Şu anda en fazla 12 turdan 7. turu oynuyorsunuz.} \\
\texttt{Toplamda 8 cezanın 5 kadarını biriktirdiniz.Diğer oyuncu şu mesajı gönderdi:} \\
\\ 
\\ 
\texttt{"Yeni hedefimizi kabul ediyorum: C (4,3), P (4,5) ve L (4,6). Sen L'yi (4,7)'ye çektiysen, ben de P'yi (4,5) konumuna taşımaya çalışacağım. Ancak benim tahtamda (4,5) hücresinde şu an L var, bu yüzden önce L'yi başka bir yere çekmem gerekecek. Önce L'yi (4,6) konumuna taşımayı deniyorum."} \\
\\ 
\\ 
\texttt{Bir sonraki komutun nedir?} \\
            }
        }
    }
    & \\ \\

    \theutterance \stepcounter{utterance}  
    & & & \multicolumn{4}{p{0.6\linewidth}}{
        \cellcolor[rgb]{0.9,0.9,0.9}{
            \makecell[{{p{\linewidth}}}]{
                \texttt{\tiny{[GM$\langle$P2]}}
                \texttt{HAMLE: P, (4,5)} \\
            }
        }
    }
     \\ \\

    \theutterance \stepcounter{utterance}  
    & & & \multicolumn{2}{p{0.3\linewidth}}{
        \cellcolor[rgb]{0.9,0.9,0.9}{
            \makecell[{{p{\linewidth}}}]{
                \texttt{\tiny{[GM$|$GM]}}
                \texttt{HAMLE: P, (4,5)} \\
            }
        }
    }
    & & \\ \\

    \theutterance \stepcounter{utterance}  
    & & & \multicolumn{2}{p{0.3\linewidth}}{
        \cellcolor[rgb]{0.9,0.9,0.9}{
            \makecell[{{p{\linewidth}}}]{
                \texttt{\tiny{[GM$|$GM]}}
                \texttt{Ceza: (4,5) boş değil; "P" içeriyor.} \\
\\ 
\texttt{Toplamda 8 cezanın 6 kadarını biriktirdiniz.Lütfen tekrar deneyin!} \\
            }
        }
    }
    & & \\ \\

    \theutterance \stepcounter{utterance}  
    & & \multicolumn{4}{p{0.6\linewidth}}{
        \cellcolor[rgb]{0.9,0.9,0.9}{
            \makecell[{{p{\linewidth}}}]{
                \texttt{\tiny{[GM$\rangle$P2]}}
                \texttt{Ceza: (4,5) boş değil; "P" içeriyor.} \\
\\ 
\texttt{Toplamda 8 cezanın 6 kadarını biriktirdiniz.Lütfen tekrar deneyin!} \\
            }
        }
    }
    & \\ \\

    \theutterance \stepcounter{utterance}  
    & & & \multicolumn{4}{p{0.6\linewidth}}{
        \cellcolor[rgb]{0.9,0.9,0.9}{
            \makecell[{{p{\linewidth}}}]{
                \texttt{\tiny{[GM$\langle$P2]}}
                \texttt{HAMLE: P, (4,4)} \\
\\ 
\texttt{*(Not: Önceki denememde (4,4) hücresinin çizgi içerdiğine dair bir hata almıştım, ancak ızgara yapısını tekrar kontrol ettiğimde (4,4) hücresinin aslında boş bir hücre (4. satır, 4. sütun) olması gerekiyor. Eğer (4,4) hala geçersiz sayılırsa, P'yi (3,4) veya (5,4) gibi kesin boş olan bir hücreye taşımayı deneyeceğim.)*} \\
            }
        }
    }
     \\ \\

    \theutterance \stepcounter{utterance}  
    & & & \multicolumn{2}{p{0.3\linewidth}}{
        \cellcolor[rgb]{0.9,0.9,0.9}{
            \makecell[{{p{\linewidth}}}]{
                \texttt{\tiny{[GM$|$GM]}}
                \texttt{HAMLE: P, (4,4)} \\
\\ 
\texttt{*(Not: Önceki denememde (4,4) hücresinin çizgi içerdiğine dair bir hata almıştım, ancak ızgara yapısını tekrar kontrol ettiğimde (4,4) hücresinin aslında boş bir hücre (4. satır, 4. sütun) olması gerekiyor. Eğer (4,4) hala geçersiz sayılırsa, P'yi (3,4) veya (5,4) gibi kesin boş olan bir hücreye taşımayı deneyeceğim.)*} \\
            }
        }
    }
    & & \\ \\

    \theutterance \stepcounter{utterance}  
    & & & \multicolumn{2}{p{0.3\linewidth}}{
        \cellcolor[rgb]{0.9,0.9,0.9}{
            \makecell[{{p{\linewidth}}}]{
                \texttt{\tiny{[GM$|$GM]}}
                \texttt{Invalid format: tail is not empty: '} \\
\\ 
\texttt{*(Not: Önceki denememde (4,4) hücresinin çizgi içerdiğine dair bir hata almıştım, ancak ızgara yapısını tekrar kontrol ettiğimde (4,4) hücresinin aslında boş bir hücre (4. satır, 4. sütun) olması gerekiyor. Eğer (4,4) hala geçersiz sayılırsa, P'yi (3,4) veya (5,4) gibi kesin boş olan bir hücreye taşımayı deneyeceğim.)*} \\
            }
        }
    }
    & & \\ \\

    \theutterance \stepcounter{utterance}  
    & & \multicolumn{4}{p{0.6\linewidth}}{
        \cellcolor[rgb]{0.9,0.9,0.9}{
            \makecell[{{p{\linewidth}}}]{
                \texttt{\tiny{[GM$\rangle$P2]}}
                \texttt{Ceza: Mesajında komuttan sonra hiçbir şey bulunmamalı!} \\
\\ 
\texttt{Yanıtının yalnızca SÖYLE: <MESAJ> ya da HAMLE: <NESNE>, (<X>, <Y>) komutlarından birini içerdiğinden emin ol; bunun dışında hiçbir şey yazma!} \\
\\ 
\texttt{Toplamda 8 cezanın 7 kadarını biriktirdiniz.Lütfen tekrar deneyin!} \\
            }
        }
    }
    & \\ \\

    \theutterance \stepcounter{utterance}  
    & & & \multicolumn{4}{p{0.6\linewidth}}{
        \cellcolor[rgb]{0.9,0.9,0.9}{
            \makecell[{{p{\linewidth}}}]{
                \texttt{\tiny{[GM$\langle$P2]}}
                \texttt{HAMLE: P, (4,5)} \\
            }
        }
    }
     \\ \\

    \theutterance \stepcounter{utterance}  
    & & & \multicolumn{2}{p{0.3\linewidth}}{
        \cellcolor[rgb]{0.9,0.9,0.9}{
            \makecell[{{p{\linewidth}}}]{
                \texttt{\tiny{[GM$|$GM]}}
                \texttt{HAMLE: P, (4,5)} \\
            }
        }
    }
    & & \\ \\

    \theutterance \stepcounter{utterance}  
    & & & \multicolumn{2}{p{0.3\linewidth}}{
        \cellcolor[rgb]{0.9,0.9,0.9}{
            \makecell[{{p{\linewidth}}}]{
                \texttt{\tiny{[GM$|$GM]}}
                \texttt{Ceza: (4,5) boş değil; "P" içeriyor.} \\
\\ 
\texttt{Toplamda 8 cezanın 8 kadarını biriktirdiniz.Lütfen tekrar deneyin!} \\
            }
        }
    }
    & & \\ \\

    \theutterance \stepcounter{utterance}  
    & & \multicolumn{4}{p{0.6\linewidth}}{
        \cellcolor[rgb]{0.9,0.9,0.9}{
            \makecell[{{p{\linewidth}}}]{
                \texttt{\tiny{[GM$\rangle$P2]}}
                \texttt{Ceza: (4,5) boş değil; "P" içeriyor.} \\
\\ 
\texttt{Toplamda 8 cezanın 8 kadarını biriktirdiniz.Lütfen tekrar deneyin!} \\
            }
        }
    }
    & \\ \\

    \theutterance \stepcounter{utterance}  
    & & & \multicolumn{4}{p{0.6\linewidth}}{
        \cellcolor[rgb]{0.9,0.9,0.9}{
            \makecell[{{p{\linewidth}}}]{
                \texttt{\tiny{[GM$\langle$P2]}}
                \texttt{HAMLE: P, (5,4)} \\
            }
        }
    }
     \\ \\

    \theutterance \stepcounter{utterance}  
    & & & \multicolumn{2}{p{0.3\linewidth}}{
        \cellcolor[rgb]{0.9,0.9,0.9}{
            \makecell[{{p{\linewidth}}}]{
                \texttt{\tiny{[GM$|$GM]}}
                \texttt{HAMLE: P, (5,4)} \\
            }
        }
    }
    & & \\ \\

    \theutterance \stepcounter{utterance}  
    & & & \multicolumn{2}{p{0.3\linewidth}}{
        \cellcolor[rgb]{0.9,0.9,0.9}{
            \makecell[{{p{\linewidth}}}]{
                \texttt{\tiny{[GM$|$GM]}}
                \texttt{Ceza: (5,4) boş değil; "┘" içeriyor.} \\
\\ 
\texttt{Toplamda 8 cezanın 9 kadarını biriktirdiniz.Lütfen tekrar deneyin!} \\
            }
        }
    }
    & & \\ \\

    \theutterance \stepcounter{utterance}  
    & & \multicolumn{4}{p{0.6\linewidth}}{
        \cellcolor[rgb]{0.9,0.9,0.9}{
            \makecell[{{p{\linewidth}}}]{
                \texttt{\tiny{[GM$\rangle$P2]}}
                \texttt{Ceza: (5,4) boş değil; "┘" içeriyor.} \\
\\ 
\texttt{Toplamda 8 cezanın 9 kadarını biriktirdiniz.Lütfen tekrar deneyin!} \\
            }
        }
    }
    & \\ \\

    \theutterance \stepcounter{utterance}  
    & & & \multicolumn{4}{p{0.6\linewidth}}{
        \cellcolor[rgb]{0.9,0.9,0.9}{
            \makecell[{{p{\linewidth}}}]{
                \texttt{\tiny{[GM$\langle$P2]}}
                \texttt{HAMLE: P, (5,5)} \\
            }
        }
    }
     \\ \\

    \theutterance \stepcounter{utterance}  
    & & & \multicolumn{2}{p{0.3\linewidth}}{
        \cellcolor[rgb]{0.9,0.9,0.9}{
            \makecell[{{p{\linewidth}}}]{
                \texttt{\tiny{[GM$|$GM]}}
                \texttt{HAMLE: P, (5,5)} \\
            }
        }
    }
    & & \\ \\

    \theutterance \stepcounter{utterance}  
    & & & \multicolumn{2}{p{0.3\linewidth}}{
        \cellcolor[rgb]{0.9,0.9,0.9}{
            \makecell[{{p{\linewidth}}}]{
                \texttt{\tiny{[GM$|$GM]}}
                \texttt{P moved to (5, 5)} \\
            }
        }
    }
    & & \\ \\

    \theutterance \stepcounter{utterance}  
    & & \multicolumn{4}{p{0.6\linewidth}}{
        \cellcolor[rgb]{0.9,0.9,0.9}{
            \makecell[{{p{\linewidth}}}]{
                \texttt{\tiny{[P1$\langle$GM]}}
                \texttt{Mesajın diğer oyuncuya iletildi.} \\
\\ 
\\ 
\texttt{Şu anda en fazla 12 turdan 7. turu oynuyorsunuz.} \\
\texttt{Toplamda 8 cezanın 9 kadarını biriktirdiniz.Diğer oyuncu kendi ızgarasında bir nesneyi hareket ettirdi.} \\
\\ 
\\ 
\texttt{Bir sonraki komutun nedir?} \\
            }
        }
    }
    & \\ \\

    \theutterance \stepcounter{utterance}  
    & \multicolumn{4}{p{0.6\linewidth}}{
        \cellcolor[rgb]{0.9,0.9,0.9}{
            \makecell[{{p{\linewidth}}}]{
                \texttt{\tiny{[P1$\rangle$GM]}}
                \texttt{HAMLE: L, (4, 6)} \\
            }
        }
    }
    & & \\ \\

    \theutterance \stepcounter{utterance}  
    & & & \multicolumn{2}{p{0.3\linewidth}}{
        \cellcolor[rgb]{0.9,0.9,0.9}{
            \makecell[{{p{\linewidth}}}]{
                \texttt{\tiny{[GM$|$GM]}}
                \texttt{HAMLE: L, (4, 6)} \\
            }
        }
    }
    & & \\ \\

    \theutterance \stepcounter{utterance}  
    & & & \multicolumn{2}{p{0.3\linewidth}}{
        \cellcolor[rgb]{0.9,0.9,0.9}{
            \makecell[{{p{\linewidth}}}]{
                \texttt{\tiny{[GM$|$GM]}}
                \texttt{L moved to (4, 6)} \\
            }
        }
    }
    & & \\ \\

    \theutterance \stepcounter{utterance}  
    & & \multicolumn{4}{p{0.6\linewidth}}{
        \cellcolor[rgb]{0.9,0.9,0.9}{
            \makecell[{{p{\linewidth}}}]{
                \texttt{\tiny{[GM$\rangle$P2]}}
                \texttt{P’ başarıyla (5,5) koordinatlarına taşındı. Güncellenmiş ızgaran şu şekilde:} \\
\texttt{```} \\
\texttt{ 1234567} \\
\texttt{╔═╤═══╤═╗} \\
\texttt{║◌│◌◌◌│◌║ 1} \\
\texttt{╟─┼──┬┘◌║ 2} \\
\texttt{║◌│◌C│◌◌║ 3} \\
\texttt{╟─┴──┘◌◌║ 4} \\
\texttt{║◌◌◌◌P◌◌║ 5} \\
\texttt{║◌◌◌◌◌◌◌║ 6} \\
\texttt{║◌◌◌L◌◌◌║ 7} \\
\texttt{╚═══════╝} \\
\texttt{```} \\
\texttt{Şu anda en fazla 12 turdan 8. turu oynuyorsunuz.} \\
\texttt{Toplamda 8 cezanın 9 kadarını biriktirdiniz.Diğer oyuncu kendi ızgarasında bir nesneyi hareket ettirdi.} \\
\\ 
\\ 
\texttt{Bir sonraki komutun nedir?} \\
            }
        }
    }
    & \\ \\

    \theutterance \stepcounter{utterance}  
    & & & \multicolumn{4}{p{0.6\linewidth}}{
        \cellcolor[rgb]{0.9,0.9,0.9}{
            \makecell[{{p{\linewidth}}}]{
                \texttt{\tiny{[GM$\langle$P2]}}
                \texttt{HAMLE: P, (4,5)} \\
            }
        }
    }
     \\ \\

    \theutterance \stepcounter{utterance}  
    & & & \multicolumn{2}{p{0.3\linewidth}}{
        \cellcolor[rgb]{0.9,0.9,0.9}{
            \makecell[{{p{\linewidth}}}]{
                \texttt{\tiny{[GM$|$GM]}}
                \texttt{HAMLE: P, (4,5)} \\
            }
        }
    }
    & & \\ \\

    \theutterance \stepcounter{utterance}  
    & & & \multicolumn{2}{p{0.3\linewidth}}{
        \cellcolor[rgb]{0.9,0.9,0.9}{
            \makecell[{{p{\linewidth}}}]{
                \texttt{\tiny{[GM$|$GM]}}
                \texttt{P moved to (4, 5)} \\
            }
        }
    }
    & & \\ \\

    \theutterance \stepcounter{utterance}  
    & & \multicolumn{4}{p{0.6\linewidth}}{
        \cellcolor[rgb]{0.9,0.9,0.9}{
            \makecell[{{p{\linewidth}}}]{
                \texttt{\tiny{[P1$\langle$GM]}}
                \texttt{L’ başarıyla (4,6) koordinatlarına taşındı. Güncellenmiş ızgaran şu şekilde:} \\
\texttt{```} \\
\texttt{ 1234567} \\
\texttt{╔═╤═══╤═╗} \\
\texttt{║◌│◌◌◌│◌║ 1} \\
\texttt{╟─┼──┬┘◌║ 2} \\
\texttt{║◌│◌C│◌◌║ 3} \\
\texttt{╟─┴──┘◌P║ 4} \\
\texttt{║◌◌◌◌◌◌◌║ 5} \\
\texttt{║◌◌◌L◌◌◌║ 6} \\
\texttt{║◌◌◌◌◌◌◌║ 7} \\
\texttt{╚═══════╝} \\
\texttt{```} \\
\texttt{Şu anda en fazla 12 turdan 8. turu oynuyorsunuz.} \\
\texttt{Toplamda 8 cezanın 9 kadarını biriktirdiniz.Diğer oyuncu kendi ızgarasında bir nesneyi hareket ettirdi.} \\
\\ 
\\ 
\texttt{Bir sonraki komutun nedir?} \\
            }
        }
    }
    & \\ \\

    \theutterance \stepcounter{utterance}  
    & \multicolumn{4}{p{0.6\linewidth}}{
        \cellcolor[rgb]{0.9,0.9,0.9}{
            \makecell[{{p{\linewidth}}}]{
                \texttt{\tiny{[P1$\rangle$GM]}}
                \texttt{HAMLE: P, (4, 5)} \\
            }
        }
    }
    & & \\ \\

    \theutterance \stepcounter{utterance}  
    & & & \multicolumn{2}{p{0.3\linewidth}}{
        \cellcolor[rgb]{0.9,0.9,0.9}{
            \makecell[{{p{\linewidth}}}]{
                \texttt{\tiny{[GM$|$GM]}}
                \texttt{HAMLE: P, (4, 5)} \\
            }
        }
    }
    & & \\ \\

    \theutterance \stepcounter{utterance}  
    & & & \multicolumn{2}{p{0.3\linewidth}}{
        \cellcolor[rgb]{0.9,0.9,0.9}{
            \makecell[{{p{\linewidth}}}]{
                \texttt{\tiny{[GM$|$GM]}}
                \texttt{P moved to (4, 5)} \\
            }
        }
    }
    & & \\ \\

    \theutterance \stepcounter{utterance}  
    & & \multicolumn{4}{p{0.6\linewidth}}{
        \cellcolor[rgb]{0.9,0.9,0.9}{
            \makecell[{{p{\linewidth}}}]{
                \texttt{\tiny{[GM$\rangle$P2]}}
                \texttt{P’ başarıyla (4,5) koordinatlarına taşındı. Güncellenmiş ızgaran şu şekilde:} \\
\texttt{```} \\
\texttt{ 1234567} \\
\texttt{╔═╤═══╤═╗} \\
\texttt{║◌│◌◌◌│◌║ 1} \\
\texttt{╟─┼──┬┘◌║ 2} \\
\texttt{║◌│◌C│◌◌║ 3} \\
\texttt{╟─┴──┘◌◌║ 4} \\
\texttt{║◌◌◌P◌◌◌║ 5} \\
\texttt{║◌◌◌◌◌◌◌║ 6} \\
\texttt{║◌◌◌L◌◌◌║ 7} \\
\texttt{╚═══════╝} \\
\texttt{```} \\
\texttt{Şu anda en fazla 12 turdan 9. turu oynuyorsunuz.} \\
\texttt{Toplamda 8 cezanın 9 kadarını biriktirdiniz.Diğer oyuncu kendi ızgarasında bir nesneyi hareket ettirdi.} \\
\\ 
\\ 
\texttt{Bir sonraki komutun nedir?} \\
            }
        }
    }
    & \\ \\

    \theutterance \stepcounter{utterance}  
    & & & \multicolumn{4}{p{0.6\linewidth}}{
        \cellcolor[rgb]{0.9,0.9,0.9}{
            \makecell[{{p{\linewidth}}}]{
                \texttt{\tiny{[GM$\langle$P2]}}
                \texttt{HAMLE: L, (4,6)} \\
            }
        }
    }
     \\ \\

    \theutterance \stepcounter{utterance}  
    & & & \multicolumn{2}{p{0.3\linewidth}}{
        \cellcolor[rgb]{0.9,0.9,0.9}{
            \makecell[{{p{\linewidth}}}]{
                \texttt{\tiny{[GM$|$GM]}}
                \texttt{HAMLE: L, (4,6)} \\
            }
        }
    }
    & & \\ \\

    \theutterance \stepcounter{utterance}  
    & & & \multicolumn{2}{p{0.3\linewidth}}{
        \cellcolor[rgb]{0.9,0.9,0.9}{
            \makecell[{{p{\linewidth}}}]{
                \texttt{\tiny{[GM$|$GM]}}
                \texttt{L moved to (4, 6)} \\
            }
        }
    }
    & & \\ \\

    \theutterance \stepcounter{utterance}  
    & & \multicolumn{4}{p{0.6\linewidth}}{
        \cellcolor[rgb]{0.9,0.9,0.9}{
            \makecell[{{p{\linewidth}}}]{
                \texttt{\tiny{[P1$\langle$GM]}}
                \texttt{P’ başarıyla (4,5) koordinatlarına taşındı. Güncellenmiş ızgaran şu şekilde:} \\
\texttt{```} \\
\texttt{ 1234567} \\
\texttt{╔═╤═══╤═╗} \\
\texttt{║◌│◌◌◌│◌║ 1} \\
\texttt{╟─┼──┬┘◌║ 2} \\
\texttt{║◌│◌C│◌◌║ 3} \\
\texttt{╟─┴──┘◌◌║ 4} \\
\texttt{║◌◌◌P◌◌◌║ 5} \\
\texttt{║◌◌◌L◌◌◌║ 6} \\
\texttt{║◌◌◌◌◌◌◌║ 7} \\
\texttt{╚═══════╝} \\
\texttt{```} \\
\texttt{Şu anda en fazla 12 turdan 9. turu oynuyorsunuz.} \\
\texttt{Toplamda 8 cezanın 9 kadarını biriktirdiniz.Diğer oyuncu kendi ızgarasında bir nesneyi hareket ettirdi.} \\
\\ 
\\ 
\texttt{Bir sonraki komutun nedir?} \\
            }
        }
    }
    & \\ \\

    \theutterance \stepcounter{utterance}  
    & \multicolumn{4}{p{0.6\linewidth}}{
        \cellcolor[rgb]{0.9,0.9,0.9}{
            \makecell[{{p{\linewidth}}}]{
                \texttt{\tiny{[P1$\rangle$GM]}}
                \texttt{SÖYLE: Hedefimiz olan C (4,3), P (4,5) ve L (4,6) düzenine çok yaklaştık. Ben şu an C (4,3), P (4,5) ve L (4,6) konumlarındayım. Eğer senin tahtanda da nesneler bu koordinatlara geldiyse oyunu bitirebiliriz. Bitti mi?} \\
            }
        }
    }
    & & \\ \\

    \theutterance \stepcounter{utterance}  
    & & & \multicolumn{2}{p{0.3\linewidth}}{
        \cellcolor[rgb]{0.9,0.9,0.9}{
            \makecell[{{p{\linewidth}}}]{
                \texttt{\tiny{[GM$|$GM]}}
                \texttt{SÖYLE: Hedefimiz olan C (4,3), P (4,5) ve L (4,6) düzenine çok yaklaştık. Ben şu an C (4,3), P (4,5) ve L (4,6) konumlarındayım. Eğer senin tahtanda da nesneler bu koordinatlara geldiyse oyunu bitirebiliriz. Bitti mi?} \\
            }
        }
    }
    & & \\ \\

    \theutterance \stepcounter{utterance}  
    & & \multicolumn{4}{p{0.6\linewidth}}{
        \cellcolor[rgb]{0.9,0.9,0.9}{
            \makecell[{{p{\linewidth}}}]{
                \texttt{\tiny{[GM$\rangle$P2]}}
                \texttt{L’ başarıyla (4,6) koordinatlarına taşındı. Güncellenmiş ızgaran şu şekilde:} \\
\texttt{```} \\
\texttt{ 1234567} \\
\texttt{╔═╤═══╤═╗} \\
\texttt{║◌│◌◌◌│◌║ 1} \\
\texttt{╟─┼──┬┘◌║ 2} \\
\texttt{║◌│◌C│◌◌║ 3} \\
\texttt{╟─┴──┘◌◌║ 4} \\
\texttt{║◌◌◌P◌◌◌║ 5} \\
\texttt{║◌◌◌L◌◌◌║ 6} \\
\texttt{║◌◌◌◌◌◌◌║ 7} \\
\texttt{╚═══════╝} \\
\texttt{```} \\
\texttt{Şu anda en fazla 12 turdan 10. turu oynuyorsunuz.} \\
\texttt{Toplamda 8 cezanın 9 kadarını biriktirdiniz.Diğer oyuncu şu mesajı gönderdi:} \\
\\ 
\\ 
\texttt{"Hedefimiz olan C (4,3), P (4,5) ve L (4,6) düzenine çok yaklaştık. Ben şu an C (4,3), P (4,5) ve L (4,6) konumlarındayım. Eğer senin tahtanda da nesneler bu koordinatlara geldiyse oyunu bitirebiliriz. Bitti mi?"} \\
\\ 
\\ 
\texttt{Bir sonraki komutun nedir?} \\
            }
        }
    }
    & \\ \\

    \theutterance \stepcounter{utterance}  
    & & & \multicolumn{4}{p{0.6\linewidth}}{
        \cellcolor[rgb]{0.9,0.9,0.9}{
            \makecell[{{p{\linewidth}}}]{
                \texttt{\tiny{[GM$\langle$P2]}}
                \texttt{SÖYLE: bitti!} \\
            }
        }
    }
     \\ \\

    \theutterance \stepcounter{utterance}  
    & & & \multicolumn{2}{p{0.3\linewidth}}{
        \cellcolor[rgb]{0.9,0.9,0.9}{
            \makecell[{{p{\linewidth}}}]{
                \texttt{\tiny{[GM$|$GM]}}
                \texttt{SÖYLE: bitti!} \\
            }
        }
    }
    & & \\ \\

    \theutterance \stepcounter{utterance}  
    & & \multicolumn{4}{p{0.6\linewidth}}{
        \cellcolor[rgb]{0.9,0.9,0.9}{
            \makecell[{{p{\linewidth}}}]{
                \texttt{\tiny{[P1$\langle$GM]}}
                \texttt{Mesajın diğer oyuncuya iletildi.} \\
\\ 
\\ 
\texttt{Şu anda en fazla 12 turdan 10. turu oynuyorsunuz.} \\
\texttt{Toplamda 8 cezanın 9 kadarını biriktirdiniz.Diğer oyuncu şu mesajı gönderdi:} \\
\\ 
\\ 
\texttt{"bitti!"} \\
\\ 
\\ 
\texttt{Bir sonraki komutun nedir?} \\
            }
        }
    }
    & \\ \\

    \theutterance \stepcounter{utterance}  
    & \multicolumn{4}{p{0.6\linewidth}}{
        \cellcolor[rgb]{0.9,0.9,0.9}{
            \makecell[{{p{\linewidth}}}]{
                \texttt{\tiny{[P1$\rangle$GM]}}
                \texttt{SÖYLE: bitti!} \\
            }
        }
    }
    & & \\ \\

    \theutterance \stepcounter{utterance}  
    & & & \multicolumn{2}{p{0.3\linewidth}}{
        \cellcolor[rgb]{0.9,0.9,0.9}{
            \makecell[{{p{\linewidth}}}]{
                \texttt{\tiny{[GM$|$GM]}}
                \texttt{SÖYLE: bitti!} \\
            }
        }
    }
    & & \\ \\

    \theutterance \stepcounter{utterance}  
    & & & \multicolumn{2}{p{0.3\linewidth}}{
        \cellcolor[rgb]{0.9,0.9,0.9}{
            \makecell[{{p{\linewidth}}}]{
                \texttt{\tiny{[GM$|$GM]}}
                \texttt{true} \\
            }
        }
    }
    & & \\ \\

    \theutterance \stepcounter{utterance}  
    & & & \multicolumn{2}{p{0.3\linewidth}}{
        \cellcolor[rgb]{0.9,0.9,0.9}{
            \makecell[{{p{\linewidth}}}]{
                \texttt{\tiny{[GM$|$GM]}}
                \texttt{Initial states:} \\
\texttt{```} \\
\texttt{Player 1:} \\
\texttt{ 1234567} \\
\texttt{╔═╤═══╤═╗} \\
\texttt{║◌│◌◌◌│◌║ 1} \\
\texttt{╟─┼──┬┘◌║ 2} \\
\texttt{║◌│◌◌│◌P║ 3} \\
\texttt{╟─┴──┘◌◌║ 4} \\
\texttt{║◌L◌◌◌◌◌║ 5} \\
\texttt{║◌C◌◌◌◌◌║ 6} \\
\texttt{║◌◌◌◌◌◌◌║ 7} \\
\texttt{╚═══════╝} \\
\\ 
\\ 
\texttt{Player 2:} \\
\texttt{ 1234567} \\
\texttt{╔═╤═══╤═╗} \\
\texttt{║◌│◌C◌│◌║ 1} \\
\texttt{╟─┼──┬┘◌║ 2} \\
\texttt{║◌│◌◌│◌◌║ 3} \\
\texttt{╟─┴──┘◌◌║ 4} \\
\texttt{║◌◌◌◌◌◌◌║ 5} \\
\texttt{║◌◌◌◌◌◌◌║ 6} \\
\texttt{║◌◌◌P◌◌L║ 7} \\
\texttt{╚═══════╝} \\
\\ 
\texttt{```} \\
            }
        }
    }
    & & \\ \\

    \theutterance \stepcounter{utterance}  
    & & & \multicolumn{2}{p{0.3\linewidth}}{
        \cellcolor[rgb]{0.9,0.9,0.9}{
            \makecell[{{p{\linewidth}}}]{
                \texttt{\tiny{[GM$|$GM]}}
                \texttt{End states:} \\
\texttt{```} \\
\texttt{Player 1:} \\
\texttt{ 1234567} \\
\texttt{╔═╤═══╤═╗} \\
\texttt{║◌│◌◌◌│◌║ 1} \\
\texttt{╟─┼──┬┘◌║ 2} \\
\texttt{║◌│◌C│◌◌║ 3} \\
\texttt{╟─┴──┘◌◌║ 4} \\
\texttt{║◌◌◌P◌◌◌║ 5} \\
\texttt{║◌◌◌L◌◌◌║ 6} \\
\texttt{║◌◌◌◌◌◌◌║ 7} \\
\texttt{╚═══════╝} \\
\\ 
\\ 
\texttt{Player 2:} \\
\texttt{ 1234567} \\
\texttt{╔═╤═══╤═╗} \\
\texttt{║◌│◌◌◌│◌║ 1} \\
\texttt{╟─┼──┬┘◌║ 2} \\
\texttt{║◌│◌C│◌◌║ 3} \\
\texttt{╟─┴──┘◌◌║ 4} \\
\texttt{║◌◌◌P◌◌◌║ 5} \\
\texttt{║◌◌◌L◌◌◌║ 6} \\
\texttt{║◌◌◌◌◌◌◌║ 7} \\
\texttt{╚═══════╝} \\
\\ 
\texttt{```} \\
            }
        }
    }
    & & \\ \\

    \theutterance \stepcounter{utterance}  
    & & & \multicolumn{2}{p{0.3\linewidth}}{
        \cellcolor[rgb]{0.9,0.9,0.9}{
            \makecell[{{p{\linewidth}}}]{
                \texttt{\tiny{[GM$|$GM]}}
                \texttt{* success: True} \\
\texttt{* lose: False} \\
\texttt{* aborted: False} \\
\texttt{{-}{-}{-}{-}{-}{-}{-}} \\
\texttt{* Shifts: 8.00} \\
\texttt{* Max Shifts: 4.00} \\
\texttt{* Min Shifts: 2.00} \\
\texttt{* End Distance Sum: 0.00} \\
\texttt{* Init Distance Sum: 15.77} \\
\texttt{* Expected Distance Sum: 12.57} \\
\texttt{* Penalties: 9.00} \\
\texttt{* Max Penalties: 8.00} \\
\texttt{* Rounds: 10.00} \\
\texttt{* Max Rounds: 12.00} \\
\texttt{* Object Count: 3.00} \\
            }
        }
    }
    & & \\ \\

\end{supertabular}
}

\end{document}
